\documentclass[11pt,a4paper]{article}
\usepackage[margin=22mm]{geometry}
\usepackage{amsmath,amssymb}
\usepackage{fontspec}
\usepackage{polyglossia}
\setdefaultlanguage{russian}
\setmainfont{DejaVu Serif}
\usepackage[unicode,hidelinks]{hyperref}

\allowdisplaybreaks[1]
\raggedright
\setlength{\parskip}{0.35em}
\newcommand{\dd}{\mathop{}\!\mathrm{d}}
\newcommand{\avg}[1]{\langle #1\rangle}
\newcommand{\ket}[1]{\lvert #1\rangle}
\newcommand{\bra}[1]{\langle #1\rvert}
\newcommand{\braket}[2]{\langle #1\mid #2\rangle}
\DeclareMathOperator{\Var}{Var}

\newcommand{\solutiontitle}[4]{%
  \hypersetup{pdftitle={Квантовая механика 1. Решения ДЗ №#1}}%
  \begin{center}
    {\small Квантовая механика 1 \enspace\textperiodcentered\enspace Семестр 1\par}
    \vspace{1ex}
    {\LARGE\bfseries Решения домашнего задания №#1\par}
    \vspace{0.5ex}
    {\large #2\par}
  \end{center}
  \noindent{\small Условия приведены в файле «#3», #4.}\par\medskip
}

\begin{document}
\solutiontitle{12}{Когерентные и сжатые состояния осциллятора}{ДЗ2026(1)\_12.pdf}{с.~1}

Используем \(Q=x/\ell,\ P=\ell p/\hbar,\ \ell=\sqrt{\hbar/(m\omega)}\),
\(a=(Q+iP)/\sqrt2,\ [a,a^\dagger]=I\) и
\(H=\hbar\omega(N+1/2)\), \(N=a^\dagger a\).

\section{Флуктуации в когерентном состоянии}

Из \(a\ket\alpha=\alpha\ket\alpha\) и
\(a\ket n=\sqrt n\ket{n-1}\) следует
\(c_{n+1}=\alpha c_n/\sqrt{n+1}\). Норма равна
\(|c_0|^2\sum_n|\alpha|^{2n}/n!=|c_0|^2e^{|\alpha|^2}\),
откуда
\[
 \ket\alpha=e^{-|\alpha|^2/2}\sum_{n\ge0}
                   \frac{\alpha^n}{\sqrt{n!}}\ket n .
\]
Фаза \(e^{-i\omega t(n+1/2)}\) даёт состояние
\(e^{-i\omega t/2}\ket{\alpha_t}\), \(\alpha_t=\alpha e^{-i\omega t}\).
В нём \(\avg{a}=\alpha_t,\ \avg{a^2}=\alpha_t^2,\ \avg{N}=|\alpha_t|^2\).
Из линейных комбинаций и квадратов \(Q,P\) следует
\begin{align*}
 \avg{Q(t)}&=\sqrt2\operatorname{Re}\alpha_t,
 &\avg{P(t)}&=\sqrt2\operatorname{Im}\alpha_t,\\
 \avg{Q^2(t)}&=|\alpha_t|^2+\tfrac12+\operatorname{Re}(\alpha_t^2),
 &\avg{P^2(t)}&=|\alpha_t|^2+\tfrac12-\operatorname{Re}(\alpha_t^2).
\end{align*}
При \(\Delta Q=Q-\avg Q I,\ \Delta P=P-\avg P I\)
средние отклонения \(\avg{\Delta Q}=\avg{\Delta P}=0\)
по определению. Используя
\(2(\operatorname{Re}\alpha_t)^2=|\alpha_t|^2+\operatorname{Re}(\alpha_t^2)\)
и аналогичное тождество с минусом для мнимой части, находим
\[
 \avg{(\Delta Q)^2}=\avg{(\Delta P)^2}=\frac12,\qquad
 \sigma_Q=\sigma_P=\frac1{\sqrt2}.
\]
Последние величины суть стандартные отклонения, если под
\(\Delta Q,\Delta P\) подразумеваются их численные значения.

\section{Дисперсия энергии}

Из \(aa^\dagger=N+I\) получаем \(N^2=(a^\dagger)^2 a^2+N\).
В когерентном состоянии \(\avg{(a^\dagger)^2a^2}
=\|a^2\ket\alpha\|^2=|\alpha|^4\).
Поэтому \(\avg{N^2}=|\alpha|^4+|\alpha|^2\),
\(\Var(N)=|\alpha|^2\), и
\[
 \avg{E}=\hbar\omega(|\alpha|^2+\tfrac12),\qquad
 \Var(E)=(\hbar\omega)^2|\alpha|^2 .
\]
При эволюции \(|\alpha_t|=|\alpha|\), поэтому моменты постоянны.

\section{Унитарность экспоненты из условия}

В формуле листка \(D(\alpha)=e^B\) показатель
\(B=\alpha(a^\dagger)^2-\alpha^*a^2\) содержит квадраты
лестничных операторов.
На общей области \(B^\dagger\) формально равен \(-B\).
Унитарность можно установить через действие на волновые функции.
Пусть \(\alpha=r e^{i\theta}\), \(r=|\alpha|\).
Унитарное \(R=e^{i\theta N/2}\) удовлетворяет
\(R a^\dagger R^\dagger=e^{i\theta/2}a^\dagger\),
\(R a R^\dagger=e^{-i\theta/2}a\), что проверяется на \(\ket n\).
Значит \(B=R\, r[(a^\dagger)^2-a^2]R^\dagger\).
Для вещественного \(r\) генератор равен
\(-ir(QP+PQ)=-r(2Q\partial_Q+1)\).
Его экспоненту задаёт группа
\[
 (U_s f)(Q)=e^{-rs} f(e^{-2rs}Q).
\]
Производная по \(s\) при нуле равна действию генератора.
Замена переменной в норме показывает \(\|U_s f\|=\|f\|\),
а \(U_{-s}\) является обратным. Следовательно
\(D(\alpha)=R U_1 R^\dagger\) унитарен и
\[
 D(\alpha)^\dagger=R U_{-1}R^\dagger=e^{-B}=D(-\alpha).
\]
При \(\alpha=0\) это единичный оператор.

Обычный оператор смещения в фазовом пространстве имеет линейный
показатель \(D_{\text{см}}(\alpha)=\exp(\alpha a^\dagger-\alpha^*a)\).
Для \(\alpha=(u+iv)/\sqrt2\) его показатель равен
\(ivQ-u\partial_Q\). Группа
\(f(Q)\mapsto e^{iv(sQ-us^2/2)}f(Q-us)\) имеет эту производную
при \(s=0\). Фазовый множитель и перенос сохраняют норму,
а замена \(s\) на \(-s\) даёт обратный оператор.
Поэтому \(D_{\text{см}}^\dagger(\alpha)=D_{\text{см}}(-\alpha)\).
Квадратичная экспонента из листка описывает сжатие.

\section{Преобразование лестничных операторов}

Для \(B\) из задачи 3 правило Лейбница даёт
\([a,B]=2\alpha a^\dagger,\ [a^\dagger,B]=2\alpha^*a\).
Производная \(e^{-sB}a e^{sB}\) есть \(e^{-sB}[a,B]e^{sB}\).
Поэтому если \(A_s=e^{-sB}a e^{sB}\) и
\(C_s=e^{-sB}a^\dagger e^{sB}\), то
\(A_s'=2\alpha C_s,\ C_s'=2\alpha^*A_s\).
Получаем \(A_s''=4r^2 A_s\) с начальными значениями
\(A_0=a,\ A'_0=2\alpha a^\dagger\).
При \(s=1\) решение и его сопряжение дают
\begin{align*}
 D(\alpha)^\dagger a D(\alpha)
    &=a\cosh(2r)+(\alpha/r)a^\dagger\sinh(2r),\\
 D(\alpha)^\dagger a^\dagger D(\alpha)
    &=a^\dagger\cosh(2r)+(\alpha^*/r)a\sinh(2r).
\end{align*}
В точке \(r=0\) берём непрерывный предел.
Для линейного показателя \(B_{\text{см}}=\alpha a^\dagger-\alpha^*a\)
коммутаторы \([a,B_{\text{см}}]=\alpha I\) и
\([a^\dagger,B_{\text{см}}]=\alpha^* I\) скалярны.
Дальнейшие коммутаторы равны нулю, так что
\(D_{\text{см}}^\dagger a D_{\text{см}}=a+\alpha I\) и
\(D_{\text{см}}^\dagger a^\dagger D_{\text{см}}=a^\dagger+\alpha^*I\).

\section{Суперпозиция вакуума и второго уровня}

\textbf{(a)} Стационарные фазы задают
\[
 \ket{\Psi(t)}=c_1e^{-i\omega t/2}\ket0+c_2e^{-5i\omega t/2}\ket2 .
\]

\textbf{(b)} В координате \(Q\) вакуумная функция
\(\varphi_0=\pi^{-1/4}e^{-Q^2/2}\) получается из \(a\varphi_0=0\).
Дважды применяя \(a^\dagger=(Q-\partial_Q)/\sqrt2\) и деля на
\(\sqrt2\), находим \(\varphi_2=(2Q^2-1)\varphi_0/\sqrt2\).
При переходе от \(\dd Q\) к \(\dd x\) множитель \(\ell^{-1/2}\)
сохраняет норму. Получаем
\[
 \Psi(x,t)=\frac{\pi^{-1/4}e^{-x^2/(2\ell^2)}}{\sqrt\ell}
     \left[c_1 e^{-i\omega t/2}
          +\frac{c_2}{\sqrt2}\left(2\frac{x^2}{\ell^2}-1\right)
                            e^{-5i\omega t/2}\right].
\]

\textbf{(c)} Обе компоненты чётны, поэтому \(\avg Q=0\).
Имеем \(\avg N=2c_2^2\) и
\(\avg{a^2}=\sqrt2\,c_1c_2 e^{-2i\omega t}\),
поскольку \(a^2\ket 2=\sqrt2\ket0\).
Из \(Q^2=N+1/2+(a^2+(a^\dagger)^2)/2\) находим
\[
 \Var(Q(t))=\tfrac12+2c_2^2+\sqrt2\,c_1c_2\cos(2\omega t).
\]

\textbf{(d)} При \(t=0\) сжатие требует
\(2c_2^2+\sqrt2\,c_1c_2<0\).
Для \(c_2\ne0\) это означает противоположные знаки коэффициентов
и \(2|c_2|<\sqrt2|c_1|\). С нормировкой получаем
\[
 c_1c_2<0,\qquad 0<|c_2|<1/\sqrt3 .
\]
При равенстве на границе дисперсия равна вакуумной.

\section{Чётная суперпозиция когерентных состояний}

\textbf{(a)} Из разложения по числам заполнения
\(\braket{\alpha}{-\alpha}=e^{-\alpha^2}\sum_{n\ge0}(-\alpha^2)^n/n!
=e^{-2\alpha^2}\). Отсюда
\[
 C=\frac{1}{\sqrt{2(1+e^{-2\alpha^2})}}.
\]

\textbf{(b)} Для любых когерентных состояний \(\beta,\gamma\)
\(\bra\beta N\ket\gamma=\beta^*\gamma\,\braket{\beta}{\gamma}\).
Два диагональных элемента дают \(2\alpha^2\), два перекрёстных
дают \(-2\alpha^2e^{-2\alpha^2}\). Поэтому
\[
 \avg{E}=\hbar\omega\left(\tfrac12+
       \alpha^2\frac{1-e^{-2\alpha^2}}{1+e^{-2\alpha^2}}\right)
          =\hbar\omega\bigl(\tfrac12+\alpha^2\tanh(\alpha^2)\bigr).
\]

\textbf{(c)} По эволюции когерентных состояний из задачи 1
\[
 \ket{\Psi(t)}=C e^{-i\omega t/2}
    \bigl(\ket{\alpha e^{-i\omega t}}+\ket{-\alpha e^{-i\omega t}}\bigr).
\]
Перекрытие обоих состояний сохраняется, и \(C\) то же самое.

\textbf{(d)} В сумме сокращаются нечётные числа заполнения.
Вакуумный коэффициент \(2C e^{-\alpha^2/2}=1/\sqrt{\cosh(\alpha^2)}\)
равен \(1+O(\alpha^4)\). Следующая амплитуда есть
\(\alpha^2/\sqrt2\) с точностью до \(O(\alpha^4)\), а хвост
с \(n\ge4\) имеет норму \(O(\alpha^4)\).
При соглашении
\[
 S(\eta)=\exp\!\left[\frac{\eta^*a^2-\eta(a^\dagger)^2}{2}\right]
\]
применение ряда к вакууму даёт
\(S(\eta)\ket0=\ket0-\eta\ket2/\sqrt2+O(|\eta|^2)\).
Сопоставляя с \(\ket\Psi=\ket0+\alpha^2\ket2/\sqrt2+O(\alpha^4)\),
находим \(\eta=-\alpha^2\). При соглашении с обратным знаком
показателя параметр равен \(+\alpha^2\).

\section{Эволюция состояния, заданного на полуоси}

В листке обе ветви кусочной функции перекрываются при \(Q<0\).
Для вопроса о попадании на левую полуось берём начальную функцию
\(\Psi(Q,0)=C f(Q)\mathbf1_{Q>0}\), где \(f\) нечётна,
не равна нулю в \(L^2\) и \(f\in L^2(\mathbb R)\).
При \(Q=0\) выбор значения не влияет на вероятности.
Тогда \(C=(\int_0^\infty |f|^2\dd Q)^{-1/2}\).
Разложение по нормированным функциям Эрмита \(\varphi_n\) даёт
\[
 b_n=C\int_0^\infty \varphi_n(q)^* f(q)\,\dd q,\qquad
 \Psi(Q,t)=\sum_{n\ge0}b_n e^{-i\omega t(n+1/2)}\varphi_n(Q).
\]
При \(t=j\pi/\omega\) фаза равна \(e^{-ij\pi/2}(-1)^{jn}\).
Так как \(\varphi_n(-Q)=(-1)^n\varphi_n(Q)\), получаем
\[
 \Psi(Q,j\pi/\omega)=e^{-ij\pi/2}\Psi((-1)^j Q,0).
\]
Следовательно при \(t=(2k+1)\pi/\omega\) вероятность
слева равна единице, а при \(t=2k\pi/\omega\) равна нулю.

Полной локализации на полуоси между этими моментами нет.
Чтобы это показать, положим \(\theta=\omega t\), \(\sin\theta\ne0\).
Формула Мелера для осциллятора даёт с точностью до постоянной
фазы
\[
 \Psi(Q,t)=\frac{e^{i Q^2\cot\theta/2}}{\sqrt{|\sin\theta|}}\,
 \widehat g_\theta(Q/\sin\theta),\qquad
 g_\theta(q)=e^{iq^2\cot\theta/2}\Psi(q,0),
\]
где \(\widehat g(u)=(2\pi)^{-1/2}\int e^{-iuq}g(q)\dd q\).
Действительно, в тождестве Мелера
\[
 \sum_{n\ge0}w^n\varphi_n(Q)\varphi_n(q)
 =\frac{\exp\!\left(
 -\frac{(1+w^2)(Q^2+q^2)-4wQq}{2(1-w^2)}
                    \right)}{\sqrt{\pi(1-w^2)}}\quad (|w|<1)
\]
подставляем \(w=\rho e^{-i\theta}\), домножаем на \(e^{-i\theta/2}\)
и берём \(\rho\uparrow1\). Квадратичные фазы дают формулу выше.
Для начальной функции с компактным носителем предел ядра
можно взять под интегралом при каждом \(Q\). Для общего \(L^2\)
формула продолжается по непрерывности, поскольку спектральный
ряд сходится в норме, а преобразование Фурье унитарно.

Для ненулевой \(g\in L^2\) с носителем \(q\ge0\)
её преобразование Фурье не может обращаться в ноль на открытом
интервале. В самом деле, при \(\operatorname{Im}\zeta<0\)
аналитическое продолжение
\[
 F(\zeta)=\frac1{\sqrt{2\pi}}\int_0^\infty g(q)e^{-i\zeta q}\dd q
    =-\frac1{2\pi i}\int_\mathbb R\frac{\widehat g(u)}{u-\zeta}\,\dd u
\]
следует из Планшереля и
\(\int_0^\infty e^{i(u-\zeta)q}\dd q=i/(u-\zeta)\).
Если \(\widehat g=0\) на интервале, последний интеграл продолжается
аналитически через его внутренние точки, так как его носитель
отделён от любой компактной части интервала, а ядро квадратично
интегрируемо на бесконечности. Граничное значение \(F\) равно
\(\widehat g\) в \(L^2\) по сходимости \(g(q)e^{-\varepsilon q}\)
к \(g(q)\). Значит продолжение равно нулю на интервале.
Теорема единственности и Планшерель тогда дали бы \(g=0\).
Применение этого факта к \(g_\theta\) показывает, что при
\(\sin\theta\ne0\) обе полуоси имеют положительные вероятности.
Первый момент с вероятностью один слева равен \(t_0=\pi/\omega\).

Если подразумевалась начальная поддержка слева, формулы применяются
со сменой \(Q\to-Q\). Тогда вероятность один слева достигается
при \(t=2k\pi/\omega\), включая \(t=0\).

\end{document}
